
\documentclass{article}
\usepackage{colortbl}
\usepackage{makecell}
\usepackage{multirow}
\usepackage{supertabular}

\begin{document}

\newcounter{utterance}

\centering \large Interaction Transcript for game `matchit\_ascii', experiment `similar\_grid\_2', episode 1 with Apertus{-}v1.5{-}8B.
\vspace{24pt}

{ \footnotesize  \setcounter{utterance}{1}
\setlength{\tabcolsep}{0pt}
\begin{supertabular}{c@{$\;$}|p{.15\linewidth}@{}p{.15\linewidth}p{.15\linewidth}p{.15\linewidth}p{.15\linewidth}p{.15\linewidth}}
    \# & $\;$A & \multicolumn{4}{c}{Game Master} & $\;\:$B\\
    \hline

    \theutterance \stepcounter{utterance}  
    & & \multicolumn{4}{p{0.6\linewidth}}{
        \cellcolor[rgb]{0.9,0.9,0.9}{
            \makecell[{{p{\linewidth}}}]{
                \texttt{\tiny{[P1$\langle$GM]}}
                \texttt{Bierzesz udział w grze polegającej na wspólnym zgadywaniu.} \\
\texttt{Celem jest ustalenie, czy ta siatka i inny obrazek, który widzę tylko ja, są takie same.} \\
\texttt{Proszę najpierw opisać swoją siatkę. Następnie ja przedstawię swój opis i będziemy mogli zadawać sobie nawzajem pytania dotyczące siatek, aby ustalić, czy są takie same. Teraz zacznij swój krótki opis siatki od "OPIS:", a następnie podaj opis. Nie dodawaj nic więcej. To jest twoja siatka:} \\
\texttt{▢ ▢ X ▢ ▢} \\
\texttt{▢ ▢ X ▢ ▢} \\
\texttt{▢ ▢ X ▢ ▢} \\
\texttt{▢ ▢ X ▢ ▢} \\
\texttt{X X X X X} \\
            }
        }
    }
    & \\ \\

    \theutterance \stepcounter{utterance}  
    & \multicolumn{4}{p{0.6\linewidth}}{
        \cellcolor[rgb]{0.9,0.9,0.9}{
            \makecell[{{p{\linewidth}}}]{
                \texttt{\tiny{[P1$\rangle$GM]}}
                \texttt{OPIS: pięcioliniowa siatka, w której na każdej z pięciu linii znajdują się po dwa znaki „▢” po lewej i prawej stronie symbolu „X” pośrodku; na ostatniej linii (pięć symboli) znajdują się wyłącznie symbole „X”.} \\
            }
        }
    }
    & & \\ \\

    \theutterance \stepcounter{utterance}  
    & & & \multicolumn{2}{p{0.3\linewidth}}{
        \cellcolor[rgb]{0.9,0.9,0.9}{
            \makecell[{{p{\linewidth}}}]{
                \texttt{\tiny{[GM$|$GM]}}
                \texttt{continue} \\
            }
        }
    }
    & & \\ \\

    \theutterance \stepcounter{utterance}  
    & & \multicolumn{4}{p{0.6\linewidth}}{
        \cellcolor[rgb]{0.9,0.9,0.9}{
            \makecell[{{p{\linewidth}}}]{
                \texttt{\tiny{[GM$\rangle$P2]}}
                \texttt{Bierzesz udział w grze polegającej na wspólnym zgadywaniu.} \\
\texttt{Celem jest ustalenie, czy ta siatka i inny obrazek, który widzę tylko ja, są takie same.} \\
\texttt{Proszę najpierw opisać swoją siatkę. Następnie ja przedstawię swój opis i będziemy mogli zadawać sobie nawzajem pytania dotyczące siatek, aby ustalić, czy są takie same. Teraz zacznij swój krótki opis siatki od "OPIS:", a następnie podaj opis. Nie dodawaj nic więcej. To jest twoja siatka:} \\
\texttt{▢ ▢ X ▢ ▢} \\
\texttt{▢ ▢ X ▢ ▢} \\
\texttt{▢ ▢ ▢ ▢ ▢} \\
\texttt{▢ ▢ X X ▢} \\
\texttt{X X X X X} \\
            }
        }
    }
    & \\ \\

    \theutterance \stepcounter{utterance}  
    & & & \multicolumn{4}{p{0.6\linewidth}}{
        \cellcolor[rgb]{0.9,0.9,0.9}{
            \makecell[{{p{\linewidth}}}]{
                \texttt{\tiny{[GM$\langle$P2]}}
                \texttt{OPIS: Siatka składa się z czterech wierszy i pięciu kolumn. Zawiera znaki: puste kwadraty (▢), „X” oraz litery „X”. W pierwszym wierszu: puste, puste, X, puste, puste. W drugim wierszu: puste, puste, X, puste, puste. W trzecim wierszu: puste, puste, puste, puste, puste. W czwartym wierszu: puste, puste, X, X, puste. Dodatkowo, pod siatką znajduje się napis: „X X X X X”.} \\
            }
        }
    }
     \\ \\

    \theutterance \stepcounter{utterance}  
    & & & \multicolumn{2}{p{0.3\linewidth}}{
        \cellcolor[rgb]{0.9,0.9,0.9}{
            \makecell[{{p{\linewidth}}}]{
                \texttt{\tiny{[GM$|$GM]}}
                \texttt{continue} \\
            }
        }
    }
    & & \\ \\

    \theutterance \stepcounter{utterance}  
    & & \multicolumn{4}{p{0.6\linewidth}}{
        \cellcolor[rgb]{0.9,0.9,0.9}{
            \makecell[{{p{\linewidth}}}]{
                \texttt{\tiny{[GM$\rangle$P2]}}
                \texttt{To jest mój OPISOPIS: pięcioliniowa siatka, w której na każdej z pięciu linii znajdują się po dwa znaki „▢” po lewej i prawej stronie symbolu „X” pośrodku; na ostatniej linii (pięć symboli) znajdują się wyłącznie symbole „X”.} \\
\texttt{Teraz zadaj pytanie, aby dowiedzieć się o nowych aspektach mojej siatki, które mogą różnić się od twojej siatki. Zacznij od "PYTANIE:" i nie dodawaj nic więcej.} \\
            }
        }
    }
    & \\ \\

    \theutterance \stepcounter{utterance}  
    & & & \multicolumn{4}{p{0.6\linewidth}}{
        \cellcolor[rgb]{0.9,0.9,0.9}{
            \makecell[{{p{\linewidth}}}]{
                \texttt{\tiny{[GM$\langle$P2]}}
                \texttt{PYTANIE: Czy symbole „▢” są rozmieszczone tak samo na każdej z pięciu linii, czy występują różnice w ich rozmieszczeniu wzdłuż linii?} \\
            }
        }
    }
     \\ \\

    \theutterance \stepcounter{utterance}  
    & & & \multicolumn{2}{p{0.3\linewidth}}{
        \cellcolor[rgb]{0.9,0.9,0.9}{
            \makecell[{{p{\linewidth}}}]{
                \texttt{\tiny{[GM$|$GM]}}
                \texttt{continue} \\
            }
        }
    }
    & & \\ \\

    \theutterance \stepcounter{utterance}  
    & & \multicolumn{4}{p{0.6\linewidth}}{
        \cellcolor[rgb]{0.9,0.9,0.9}{
            \makecell[{{p{\linewidth}}}]{
                \texttt{\tiny{[P1$\langle$GM]}}
                \texttt{To jest mój OPISOPIS: Siatka składa się z czterech wierszy i pięciu kolumn. Zawiera znaki: puste kwadraty (▢), „X” oraz litery „X”. W pierwszym wierszu: puste, puste, X, puste, puste. W drugim wierszu: puste, puste, X, puste, puste. W trzecim wierszu: puste, puste, puste, puste, puste. W czwartym wierszu: puste, puste, X, X, puste. Dodatkowo, pod siatką znajduje się napis: „X X X X X”.} \\
\texttt{PYTANIE: Czy symbole „▢” są rozmieszczone tak samo na każdej z pięciu linii, czy występują różnice w ich rozmieszczeniu wzdłuż linii?Rozpocznij odpowiedź od "ODPOWIEDŹ:" i nie dodawaj nic więcej.} \\
            }
        }
    }
    & \\ \\

    \theutterance \stepcounter{utterance}  
    & \multicolumn{4}{p{0.6\linewidth}}{
        \cellcolor[rgb]{0.9,0.9,0.9}{
            \makecell[{{p{\linewidth}}}]{
                \texttt{\tiny{[P1$\rangle$GM]}}
                \texttt{ODPOWIEDŹ: Rozmieszczenie pól „▢” różni się wzdłuż linii. W pierwszych trzech wierszach „▢” znajdują się po jednej jednostce po lewej i prawej stronie „X”, co daje wzór „▢ ▢ X ▢ ▢”. W ostatnim wierszu „▢” nie występują, a „X” są uszeregowane obok siebie.} \\
            }
        }
    }
    & & \\ \\

    \theutterance \stepcounter{utterance}  
    & & & \multicolumn{2}{p{0.3\linewidth}}{
        \cellcolor[rgb]{0.9,0.9,0.9}{
            \makecell[{{p{\linewidth}}}]{
                \texttt{\tiny{[GM$|$GM]}}
                \texttt{continue} \\
            }
        }
    }
    & & \\ \\

    \theutterance \stepcounter{utterance}  
    & & \multicolumn{4}{p{0.6\linewidth}}{
        \cellcolor[rgb]{0.9,0.9,0.9}{
            \makecell[{{p{\linewidth}}}]{
                \texttt{\tiny{[P1$\langle$GM]}}
                \texttt{Teraz zadaj pytanie, aby dowiedzieć się o nowych aspektach mojej siatki, które mogą różnić się od twojej siatki. Zacznij od "PYTANIE:" i nie dodawaj nic więcej.} \\
            }
        }
    }
    & \\ \\

    \theutterance \stepcounter{utterance}  
    & \multicolumn{4}{p{0.6\linewidth}}{
        \cellcolor[rgb]{0.9,0.9,0.9}{
            \makecell[{{p{\linewidth}}}]{
                \texttt{\tiny{[P1$\rangle$GM]}}
                \texttt{PYTANIE: Czy istnieją symbole lub elementy poza pólkami „▢” i „X”, takie jak litery lub inne znaki, znajdujące się w Twojej siatce?} \\
            }
        }
    }
    & & \\ \\

    \theutterance \stepcounter{utterance}  
    & & & \multicolumn{2}{p{0.3\linewidth}}{
        \cellcolor[rgb]{0.9,0.9,0.9}{
            \makecell[{{p{\linewidth}}}]{
                \texttt{\tiny{[GM$|$GM]}}
                \texttt{continue} \\
            }
        }
    }
    & & \\ \\

    \theutterance \stepcounter{utterance}  
    & & \multicolumn{4}{p{0.6\linewidth}}{
        \cellcolor[rgb]{0.9,0.9,0.9}{
            \makecell[{{p{\linewidth}}}]{
                \texttt{\tiny{[GM$\rangle$P2]}}
                \texttt{ODPOWIEDŹ: Rozmieszczenie pól „▢” różni się wzdłuż linii. W pierwszych trzech wierszach „▢” znajdują się po jednej jednostce po lewej i prawej stronie „X”, co daje wzór „▢ ▢ X ▢ ▢”. W ostatnim wierszu „▢” nie występują, a „X” są uszeregowane obok siebie.} \\
\texttt{PYTANIE: Czy istnieją symbole lub elementy poza pólkami „▢” i „X”, takie jak litery lub inne znaki, znajdujące się w Twojej siatce?Rozpocznij odpowiedź od "ODPOWIEDŹ:" i nie dodawaj nic więcej.} \\
            }
        }
    }
    & \\ \\

    \theutterance \stepcounter{utterance}  
    & & & \multicolumn{4}{p{0.6\linewidth}}{
        \cellcolor[rgb]{0.9,0.9,0.9}{
            \makecell[{{p{\linewidth}}}]{
                \texttt{\tiny{[GM$\langle$P2]}}
                \texttt{ODPOWIEDŹ: Nie, w mojej siatce nie znajdują się żadne inne symbole ani litery poza znakami „▢” i „X”.} \\
            }
        }
    }
     \\ \\

    \theutterance \stepcounter{utterance}  
    & & & \multicolumn{2}{p{0.3\linewidth}}{
        \cellcolor[rgb]{0.9,0.9,0.9}{
            \makecell[{{p{\linewidth}}}]{
                \texttt{\tiny{[GM$|$GM]}}
                \texttt{continue} \\
            }
        }
    }
    & & \\ \\

    \theutterance \stepcounter{utterance}  
    & & \multicolumn{4}{p{0.6\linewidth}}{
        \cellcolor[rgb]{0.9,0.9,0.9}{
            \makecell[{{p{\linewidth}}}]{
                \texttt{\tiny{[GM$\rangle$P2]}}
                \texttt{Teraz zadaj pytanie, aby dowiedzieć się o nowych aspektach mojej siatki, które mogą różnić się od twojej siatki. Zacznij od "PYTANIE:" i nie dodawaj nic więcej.} \\
            }
        }
    }
    & \\ \\

    \theutterance \stepcounter{utterance}  
    & & & \multicolumn{4}{p{0.6\linewidth}}{
        \cellcolor[rgb]{0.9,0.9,0.9}{
            \makecell[{{p{\linewidth}}}]{
                \texttt{\tiny{[GM$\langle$P2]}}
                \texttt{PYTANIE: Czy w Twojej siatce istnieją poziome lub pionowe linie łączące symbole, które nie znajdują się w mojej siatce?} \\
            }
        }
    }
     \\ \\

    \theutterance \stepcounter{utterance}  
    & & & \multicolumn{2}{p{0.3\linewidth}}{
        \cellcolor[rgb]{0.9,0.9,0.9}{
            \makecell[{{p{\linewidth}}}]{
                \texttt{\tiny{[GM$|$GM]}}
                \texttt{continue} \\
            }
        }
    }
    & & \\ \\

    \theutterance \stepcounter{utterance}  
    & & \multicolumn{4}{p{0.6\linewidth}}{
        \cellcolor[rgb]{0.9,0.9,0.9}{
            \makecell[{{p{\linewidth}}}]{
                \texttt{\tiny{[P1$\langle$GM]}}
                \texttt{ODPOWIEDŹ: Nie, w mojej siatce nie znajdują się żadne inne symbole ani litery poza znakami „▢” i „X”.} \\
\texttt{PYTANIE: Czy w Twojej siatce istnieją poziome lub pionowe linie łączące symbole, które nie znajdują się w mojej siatce?Rozpocznij odpowiedź od "ODPOWIEDŹ:" i nie dodawaj nic więcej.} \\
            }
        }
    }
    & \\ \\

    \theutterance \stepcounter{utterance}  
    & \multicolumn{4}{p{0.6\linewidth}}{
        \cellcolor[rgb]{0.9,0.9,0.9}{
            \makecell[{{p{\linewidth}}}]{
                \texttt{\tiny{[P1$\rangle$GM]}}
                \texttt{ODPOWIEDŹ: Nie, w mojej siatce nie występują żadne linie łączące symbole ani dodatkowe elementy wizualne poza samymi znakami „▢” i „X”.} \\
            }
        }
    }
    & & \\ \\

    \theutterance \stepcounter{utterance}  
    & & & \multicolumn{2}{p{0.3\linewidth}}{
        \cellcolor[rgb]{0.9,0.9,0.9}{
            \makecell[{{p{\linewidth}}}]{
                \texttt{\tiny{[GM$|$GM]}}
                \texttt{continue} \\
            }
        }
    }
    & & \\ \\

    \theutterance \stepcounter{utterance}  
    & & \multicolumn{4}{p{0.6\linewidth}}{
        \cellcolor[rgb]{0.9,0.9,0.9}{
            \makecell[{{p{\linewidth}}}]{
                \texttt{\tiny{[P1$\langle$GM]}}
                \texttt{Teraz zadaj pytanie, aby dowiedzieć się o nowych aspektach mojej siatki, które mogą różnić się od twojej siatki. Zacznij od "PYTANIE:" i nie dodawaj nic więcej.} \\
            }
        }
    }
    & \\ \\

    \theutterance \stepcounter{utterance}  
    & \multicolumn{4}{p{0.6\linewidth}}{
        \cellcolor[rgb]{0.9,0.9,0.9}{
            \makecell[{{p{\linewidth}}}]{
                \texttt{\tiny{[P1$\rangle$GM]}}
                \texttt{PYTANIE: Czy moja siatka posiada klucz umieszczony obok lub poniżej niej, zawierający informacje, które pomagają zrozumieć układ znaków w siatce?} \\
            }
        }
    }
    & & \\ \\

    \theutterance \stepcounter{utterance}  
    & & & \multicolumn{2}{p{0.3\linewidth}}{
        \cellcolor[rgb]{0.9,0.9,0.9}{
            \makecell[{{p{\linewidth}}}]{
                \texttt{\tiny{[GM$|$GM]}}
                \texttt{continue} \\
            }
        }
    }
    & & \\ \\

    \theutterance \stepcounter{utterance}  
    & & \multicolumn{4}{p{0.6\linewidth}}{
        \cellcolor[rgb]{0.9,0.9,0.9}{
            \makecell[{{p{\linewidth}}}]{
                \texttt{\tiny{[GM$\rangle$P2]}}
                \texttt{ODPOWIEDŹ: Nie, w mojej siatce nie występują żadne linie łączące symbole ani dodatkowe elementy wizualne poza samymi znakami „▢” i „X”.} \\
\texttt{PYTANIE: Czy moja siatka posiada klucz umieszczony obok lub poniżej niej, zawierający informacje, które pomagają zrozumieć układ znaków w siatce?Rozpocznij odpowiedź od "ODPOWIEDŹ:" i nie dodawaj nic więcej.} \\
            }
        }
    }
    & \\ \\

    \theutterance \stepcounter{utterance}  
    & & & \multicolumn{4}{p{0.6\linewidth}}{
        \cellcolor[rgb]{0.9,0.9,0.9}{
            \makecell[{{p{\linewidth}}}]{
                \texttt{\tiny{[GM$\langle$P2]}}
                \texttt{ODPOWIEDŹ: Nie, w mojej siatce nie występuje żaden klucz ani dodatkowe informacje obok lub poniżej siatki.} \\
            }
        }
    }
     \\ \\

    \theutterance \stepcounter{utterance}  
    & & & \multicolumn{2}{p{0.3\linewidth}}{
        \cellcolor[rgb]{0.9,0.9,0.9}{
            \makecell[{{p{\linewidth}}}]{
                \texttt{\tiny{[GM$|$GM]}}
                \texttt{continue} \\
            }
        }
    }
    & & \\ \\

    \theutterance \stepcounter{utterance}  
    & & \multicolumn{4}{p{0.6\linewidth}}{
        \cellcolor[rgb]{0.9,0.9,0.9}{
            \makecell[{{p{\linewidth}}}]{
                \texttt{\tiny{[GM$\rangle$P2]}}
                \texttt{Teraz zadaj pytanie, aby dowiedzieć się o nowych aspektach mojej siatki, które mogą różnić się od twojej siatki. Zacznij od "PYTANIE:" i nie dodawaj nic więcej.} \\
            }
        }
    }
    & \\ \\

    \theutterance \stepcounter{utterance}  
    & & & \multicolumn{4}{p{0.6\linewidth}}{
        \cellcolor[rgb]{0.9,0.9,0.9}{
            \makecell[{{p{\linewidth}}}]{
                \texttt{\tiny{[GM$\langle$P2]}}
                \texttt{PYTANIE: Czy w Twojej siatce występuje bilans lub liczby pod każdą linią siatki, wskazujące na ilość symboli „X” w każdej linii?} \\
            }
        }
    }
     \\ \\

    \theutterance \stepcounter{utterance}  
    & & & \multicolumn{2}{p{0.3\linewidth}}{
        \cellcolor[rgb]{0.9,0.9,0.9}{
            \makecell[{{p{\linewidth}}}]{
                \texttt{\tiny{[GM$|$GM]}}
                \texttt{continue} \\
            }
        }
    }
    & & \\ \\

    \theutterance \stepcounter{utterance}  
    & & \multicolumn{4}{p{0.6\linewidth}}{
        \cellcolor[rgb]{0.9,0.9,0.9}{
            \makecell[{{p{\linewidth}}}]{
                \texttt{\tiny{[P1$\langle$GM]}}
                \texttt{ODPOWIEDŹ: Nie, w mojej siatce nie występuje żaden klucz ani dodatkowe informacje obok lub poniżej siatki.} \\
\texttt{PYTANIE: Czy w Twojej siatce występuje bilans lub liczby pod każdą linią siatki, wskazujące na ilość symboli „X” w każdej linii?Rozpocznij odpowiedź od "ODPOWIEDŹ:" i nie dodawaj nic więcej.} \\
            }
        }
    }
    & \\ \\

    \theutterance \stepcounter{utterance}  
    & \multicolumn{4}{p{0.6\linewidth}}{
        \cellcolor[rgb]{0.9,0.9,0.9}{
            \makecell[{{p{\linewidth}}}]{
                \texttt{\tiny{[P1$\rangle$GM]}}
                \texttt{ODPOWIEDŹ: Nie, pod każdą linią mojej siatki nie znajdują się żadne bilanse ani liczby wskazujące na ilość symboli „X”.} \\
            }
        }
    }
    & & \\ \\

    \theutterance \stepcounter{utterance}  
    & & & \multicolumn{2}{p{0.3\linewidth}}{
        \cellcolor[rgb]{0.9,0.9,0.9}{
            \makecell[{{p{\linewidth}}}]{
                \texttt{\tiny{[GM$|$GM]}}
                \texttt{continue} \\
            }
        }
    }
    & & \\ \\

    \theutterance \stepcounter{utterance}  
    & & \multicolumn{4}{p{0.6\linewidth}}{
        \cellcolor[rgb]{0.9,0.9,0.9}{
            \makecell[{{p{\linewidth}}}]{
                \texttt{\tiny{[P1$\langle$GM]}}
                \texttt{Teraz zadaj pytanie, aby dowiedzieć się o nowych aspektach mojej siatki, które mogą różnić się od twojej siatki. Zacznij od "PYTANIE:" i nie dodawaj nic więcej.} \\
            }
        }
    }
    & \\ \\

    \theutterance \stepcounter{utterance}  
    & \multicolumn{4}{p{0.6\linewidth}}{
        \cellcolor[rgb]{0.9,0.9,0.9}{
            \makecell[{{p{\linewidth}}}]{
                \texttt{\tiny{[P1$\rangle$GM]}}
                \texttt{PYTANIE: Czy moja siatka posiada oznaczenia takie jak numery linii lub kolumn, które pomagają nazwać lub zidentyfikować poszczególne elementy w siatce?} \\
            }
        }
    }
    & & \\ \\

    \theutterance \stepcounter{utterance}  
    & & & \multicolumn{2}{p{0.3\linewidth}}{
        \cellcolor[rgb]{0.9,0.9,0.9}{
            \makecell[{{p{\linewidth}}}]{
                \texttt{\tiny{[GM$|$GM]}}
                \texttt{continue} \\
            }
        }
    }
    & & \\ \\

    \theutterance \stepcounter{utterance}  
    & & \multicolumn{4}{p{0.6\linewidth}}{
        \cellcolor[rgb]{0.9,0.9,0.9}{
            \makecell[{{p{\linewidth}}}]{
                \texttt{\tiny{[GM$\rangle$P2]}}
                \texttt{ODPOWIEDŹ: Nie, pod każdą linią mojej siatki nie znajdują się żadne bilanse ani liczby wskazujące na ilość symboli „X”.} \\
\texttt{PYTANIE: Czy moja siatka posiada oznaczenia takie jak numery linii lub kolumn, które pomagają nazwać lub zidentyfikować poszczególne elementy w siatce?Rozpocznij odpowiedź od "ODPOWIEDŹ:" i nie dodawaj nic więcej.} \\
            }
        }
    }
    & \\ \\

    \theutterance \stepcounter{utterance}  
    & & & \multicolumn{4}{p{0.6\linewidth}}{
        \cellcolor[rgb]{0.9,0.9,0.9}{
            \makecell[{{p{\linewidth}}}]{
                \texttt{\tiny{[GM$\langle$P2]}}
                \texttt{ODPOWIEDŹ: Nie, moja siatka nie posiada żadnych oznaczeń takich jak numery linii lub kolumn, które mogłyby pomóc w identyfikacji poszczególnych elementów.} \\
            }
        }
    }
     \\ \\

    \theutterance \stepcounter{utterance}  
    & & & \multicolumn{2}{p{0.3\linewidth}}{
        \cellcolor[rgb]{0.9,0.9,0.9}{
            \makecell[{{p{\linewidth}}}]{
                \texttt{\tiny{[GM$|$GM]}}
                \texttt{continue} \\
            }
        }
    }
    & & \\ \\

    \theutterance \stepcounter{utterance}  
    & & \multicolumn{4}{p{0.6\linewidth}}{
        \cellcolor[rgb]{0.9,0.9,0.9}{
            \makecell[{{p{\linewidth}}}]{
                \texttt{\tiny{[GM$\rangle$P2]}}
                \texttt{Teraz podejmij decyzję. Jak myślisz: czy Twoja siatka i druga opisana siatka są takie same? Napisz "DECYZJA: ta sama siatka" jeśli uważasz, że to ta sama siatka, albo "DECYZJA: różne siatki" jeśli uważasz, że to różne siatki. Nie dodawaj niczego więcej.} \\
            }
        }
    }
    & \\ \\

    \theutterance \stepcounter{utterance}  
    & & & \multicolumn{4}{p{0.6\linewidth}}{
        \cellcolor[rgb]{0.9,0.9,0.9}{
            \makecell[{{p{\linewidth}}}]{
                \texttt{\tiny{[GM$\langle$P2]}}
                \texttt{DECYZJA: różne siatki} \\
            }
        }
    }
     \\ \\

    \theutterance \stepcounter{utterance}  
    & & & \multicolumn{2}{p{0.3\linewidth}}{
        \cellcolor[rgb]{0.9,0.9,0.9}{
            \makecell[{{p{\linewidth}}}]{
                \texttt{\tiny{[GM$|$GM]}}
                \texttt{continue} \\
            }
        }
    }
    & & \\ \\

    \theutterance \stepcounter{utterance}  
    & & & \multicolumn{2}{p{0.3\linewidth}}{
        \cellcolor[rgb]{0.9,0.9,0.9}{
            \makecell[{{p{\linewidth}}}]{
                \texttt{\tiny{[GM$|$GM]}}
                \texttt{success} \\
            }
        }
    }
    & & \\ \\

    \theutterance \stepcounter{utterance}  
    & & \multicolumn{4}{p{0.6\linewidth}}{
        \cellcolor[rgb]{0.9,0.9,0.9}{
            \makecell[{{p{\linewidth}}}]{
                \texttt{\tiny{[P1$\langle$GM]}}
                \texttt{ODPOWIEDŹ: Nie, moja siatka nie posiada żadnych oznaczeń takich jak numery linii lub kolumn, które mogłyby pomóc w identyfikacji poszczególnych elementów.} \\
\texttt{Teraz podejmij decyzję. Jak myślisz: czy Twoja siatka i druga opisana siatka są takie same? Napisz "DECYZJA: ta sama siatka" jeśli uważasz, że to ta sama siatka, albo "DECYZJA: różne siatki" jeśli uważasz, że to różne siatki. Nie dodawaj niczego więcej.} \\
            }
        }
    }
    & \\ \\

    \theutterance \stepcounter{utterance}  
    & \multicolumn{4}{p{0.6\linewidth}}{
        \cellcolor[rgb]{0.9,0.9,0.9}{
            \makecell[{{p{\linewidth}}}]{
                \texttt{\tiny{[P1$\rangle$GM]}}
                \texttt{DECYZJA: różne siatki} \\
            }
        }
    }
    & & \\ \\

    \theutterance \stepcounter{utterance}  
    & & & \multicolumn{2}{p{0.3\linewidth}}{
        \cellcolor[rgb]{0.9,0.9,0.9}{
            \makecell[{{p{\linewidth}}}]{
                \texttt{\tiny{[GM$|$GM]}}
                \texttt{continue} \\
            }
        }
    }
    & & \\ \\

    \theutterance \stepcounter{utterance}  
    & & & \multicolumn{2}{p{0.3\linewidth}}{
        \cellcolor[rgb]{0.9,0.9,0.9}{
            \makecell[{{p{\linewidth}}}]{
                \texttt{\tiny{[GM$|$GM]}}
                \texttt{success} \\
            }
        }
    }
    & & \\ \\

\end{supertabular}
}

\end{document}
